\section{The spectrum of the conic surgery models}\label{sec:spectrum}

In this section $R$, $A$ and the unit $u$ are fixed, and we study the
models $(\overline X_{\epsilon,u},\bar g_{\epsilon,u}^{R,A})$ of
Definition~\ref{def:model} as $\epsilon\to0$. We abbreviate them
$(\overline X_\epsilon,\bar g_\epsilon)$ and write $b=\dim V$, $s=\dim K$,
$m=b-s$, $c=\pi$ for the constant of \eqref{eq:link-constant}, and
$\ell_\epsilon=\log(1/\epsilon)$. The arguments follow
\cite[Sections~4--6]{wuebben-conifold}, where the cone is the Ricci-flat
conifold; we give them in full, since the cone, its link and the cap are
different here and only the properties recorded in
Section~\ref{sec:node} are used.

\subsection{The energy of a period across a cap}
For a smooth two-form $\psi$ on a region of a cone or of a cap that
contains the section $\{r=T\}$, let
\begin{equation}\label{eq:boundary-coefficient}
 \mathsf B_\psi(T)=c^{-1}\int_{L}\psi_T\wedge*_h\vartheta ,
\end{equation}
where $\psi_T$ is the pull-back of $\psi$ to the link at radius $T$,
regarded as a form on $L$. Since $*_h\vartheta$ is closed, $\mathsf
B_\psi(T)$ is the coefficient of $\vartheta$ in the harmonic part of
$\psi_T$, and it depends only on the class of $\psi_T$ when $\psi_T$ is
closed.

\begin{lemma}\label{lem:cap-energy}
Let $(V_\sigma,g_V^A)$ be a cap, with radial function $r$ on its conical
part $\{r\ge2A\}$. There is $C_A$ such that for every $T>2A$ and every
smooth two-form $\psi$ on $V_\sigma\cap\{r\le T\}$,
\begin{equation}\label{eq:cap-energy}
 \int_{\{r\le T\}}\bigl(|d\psi|^2+|d^*\psi|^2\bigr)\,dV
 \ge\frac{c\,|\mathsf B_\psi(T)|^2}{\log T+C_A}.
\end{equation}
The left side is unchanged when the metric is multiplied by a constant.
\end{lemma}
\begin{proof}
By \eqref{eq:cap-topology} there is a closed three-form $\Xi$ on the
compact set $\{r\le2A\}$ whose restriction to the boundary is
$*_h\vartheta$. Stokes' theorem and the Cauchy--Schwarz inequality give
\[
 c\,\mathsf B_\psi(2A)=\int_{\{r\le2A\}}d\psi\wedge\Xi,\qquad
 \int_{\{r\le2A\}}|d\psi|^2\ge\kappa_A|\mathsf B_\psi(2A)|^2,
 \quad\kappa_A=c^2/\|\Xi\|^2_{L^2}.
\]
On the conical part write $\psi=\psi_r+dr\wedge\beta(r)$. The component
of $d\psi$ along $dr\wedge\vartheta$ is $\mathsf B_\psi'(r)\,dr\wedge
\vartheta$, because $d_L^*\vartheta=0$. Since
$|dr\wedge\vartheta|^2=r^{-4}|\vartheta|_h^2$ and the volume density is
$r^5dr\,dV_h$,
\[
 \int_{\{2A\le r\le T\}}|d\psi|^2\ge c\int_{2A}^Tr\,|\mathsf B_\psi'(r)|^2dr .
\]
Applying the Cauchy--Schwarz inequality to
$\mathsf B_\psi(T)=\mathsf B_\psi(2A)+\int_{2A}^T\mathsf B_\psi'$ gives
\[
 |\mathsf B_\psi(T)|^2\le\Bigl(\kappa_A^{-1}+c^{-1}\log\frac T{2A}\Bigr)
 \int_{\{r\le T\}}|d\psi|^2 ,
\]
which is \eqref{eq:cap-energy}. In dimension six the integrals of
$|d\psi|^2$ and $|d^*\psi|^2$ for a two-form scale by
$\epsilon^{6-6}$ and $\epsilon^{6-4-2}$ under $g\mapsto\epsilon^2g$.
\end{proof}

\subsection{Central harmonic forms}
\begin{lemma}\label{lem:central-kernels}
For the metric $g_0^R$ on $M_0$:
\begin{enumerate}
\item every smooth square-integrable closed and coclosed one-form
vanishes;
\item the smooth square-integrable closed and coclosed two-forms form a
space $\mathcal H_R$ that maps isomorphically onto $V$; write
$\eta^R_D$ for the element with class $D$;
\item on the conical part $\{r<R\}$ of the end at $p_i$,
\begin{equation}\label{eq:cone-expansion}
 \eta^R_D=\per_i(D)\,\vartheta_i+d\beta_{D},\qquad
 |\beta_D|=O(r^{\kappa-1}),\quad|d\beta_D|=O(r^{\kappa-2}),
\end{equation}
for some $\kappa>0$, with $\beta_D$ a one-form without $dr$-component whose
restriction to each link section is coclosed, and with differentiated
bounds;
\item these forms lie in the minimal domain of $d+d^*$, that is, they are
limits in the graph norm of compactly supported smooth forms.
\end{enumerate}
\end{lemma}
\begin{proof}
On an exactly conical end write a closed and coclosed two-form as
$\psi=a(r)+dr\wedge b(r)$ and decompose
$a(r)=c_0(r)\vartheta+d_Lv(r)$ with $v(r)$ coclosed on the link; there are
no harmonic one-forms on $L$. Closedness gives $c_0'=0$ and that $b-v'$
is closed; coclosedness then makes $b-v'$ harmonic, hence zero, and
\[
 b=v',\qquad r^{-2}\Delta_Lv-v''-3r^{-1}v'=0 .
\]
On the eigenspace of coclosed one-forms on $L$ with eigenvalue
$\lambda_L>0$ the solutions are $r^{k_\pm}$ with
$k_\pm=-1\pm\sqrt{1+\lambda_L}$. The corresponding two-forms have
pointwise size $r^{k_\pm-2}$, and are square integrable near the vertex
exactly when $k_\pm>-1$, since their squared radial norm is proportional
to $\int_0^{r_0}r^{2k_\pm+1}\,dr$; this excludes $k_-<-2$. This gives
\eqref{eq:cone-expansion} with $\beta_D$ the sum of the terms with
$k_+>0$ and $\kappa=\min k_+$; convergence and the differentiated bounds
follow from the expansion in eigenforms and elliptic estimates on
annuli of fixed ratio. The term $\per_i(D)\vartheta_i$ is identified by its
class.

For (4), cut off the term $\vartheta_i$ with a function $\chi_\rho$ equal
to zero for $r\le\rho^2$ and one for $r\ge\rho$, with
$|d\chi_\rho|\le C/(r|\log\rho|)$; since $\vartheta_i$ is closed and
coclosed with $|\vartheta_i|=O(r^{-2})$, the graph-norm error is
$O(|\log\rho|^{-1/2})$. Cut off $d\beta_D$ with an ordinary cutoff at
scale $\rho$; the error is $O(\rho^{\kappa})$.

For (2), injectivity is Lemma~\ref{lem:uniqueness}. For surjectivity
consider the closure $d_{\min}$ of $d$ on compactly supported forms. The
orthogonal complement in $\ker d_{\min}$ of the closure of the image of
$d_{\min}$ consists of closed and coclosed forms, which are smooth by
elliptic regularity. The forms $\eta_D$ of
Proposition~\ref{prop:harmonic} are closed, square integrable for
$g_0^R$, which is uniformly equivalent to $g_0$, and lie in the domain of
$d_{\min}$ by the cutoff argument just given. Their images in the
quotient $\ker d_{\min}/\overline{\im d_{\min}}$ are independent, because
the pairings with compactly supported closed four-forms representing a
basis of $H^4_c(M_0)$ are continuous in $L^2$ and vanish on
$\overline{\im d_{\min}}$. Hence $\dim\mathcal H_R\ge b$.

For (1), the same analysis on an end shows that a closed and coclosed
square-integrable one-form is, near the vertex, the differential of a
sum of homogeneous harmonic functions with exponents
$-2+\sqrt{4+\lambda_L}>0$, for eigenvalues $\lambda_L$ of the link; the remaining solutions $r^{-5}dr$ and those with
negative exponents are not square integrable. Since $H^1(M_0)=0$ the
form has a global primitive, bounded near the nodes, and the cutoff
argument of Lemma~\ref{lem:uniqueness} shows that it vanishes.
\end{proof}

Let $G_R(D,E)=\langle\eta^R_D,\eta^R_E\rangle_{L^2(g_0^R)}$ and let
$\cM_R$ be its quotient on $\cQ$. Each $\eta^R_D$ minimizes the $L^2$
norm of $g_0^R$ among closed forms in the domain of $d_{\min}$ with
class $D$, and the same holds for $\eta_D$ and $g_0$. Since the two
norms differ by a factor $1+O(R)$ by \eqref{eq:interpolation-errors},
\begin{equation}\label{eq:mass-limit}
 G_R=(1+O(R))\,G,\qquad\cM_R=(1+O(R))\,\cM,
\end{equation}
as quadratic forms, and the eigenvalues $\mu_j(R)$ of $\cE$ relative to
$\cM_R$ converge to $\mu_j$ as $R\to0$.

\subsection{The caps}
\begin{lemma}\label{lem:cap-kernels}
For the complete metric $g_V^A$ on $V_\sigma$:
\begin{enumerate}
\item there are no nonzero square-integrable harmonic forms;
\item if $\psi$ is a smooth solution of $(d+d^*)\psi=0$ with
$|\psi|=O(r^{-3})$ at infinity, then the part of $\psi$ that is
homogeneous of pointwise order $r^{-3}$ has form degree three.
\end{enumerate}
\end{lemma}
\begin{proof}
The metric $g_V^A$ is exactly conical outside a compact set. By
\cite[Theorem~1A]{hhm} its space of $L^2$ harmonic forms of degree $j$ is
$H^j_c(V_\sigma)$ for $j<3$, $H^j(V_\sigma)$ for $j>3$, and the image of
$H^3_c(V_\sigma)\to H^3(V_\sigma)$ for $j=3$. Since $V_\sigma\simeq
T^*S^3$, the first two vanish, and the third is zero because the map is
the Euler class of the cotangent bundle of $S^3$. This proves (1).

For (2), first $d\psi=d^*\psi=0$ separately: for a cutoff $\chi_T$ equal
to one for $r\le T$ and supported in $r\le2T$,
$\|(d+d^*)(\chi_T\psi)\|^2\le CT^{-2}\int_{T\le r\le2T}|\psi|^2\le
CT^{-2}$, and for compactly supported forms this norm is
$\|d(\chi_T\psi)\|^2+\|d^*(\chi_T\psi)\|^2$. A closed and coclosed form
of pure degree $p$, homogeneous of pointwise order $r^{-3}$ on the cone,
is $r^{p-3}a+dr\wedge r^{p-4}b$ with $a,b$ on the link and
\[
 d_La=0,\quad(p-3)a=d_Lb,\quad d_L^*b=0,\quad d_L^*a+(p-3)b=0 .
\]
For $p\ne3$ these give $\Delta_Lb=-(p-3)^2b$, so $b=0$ and $a=0$.
\end{proof}

\subsection{Compactness}
\begin{proposition}\label{prop:compactness}
Let $p\in\{1,2\}$ and let $\psi^1_\epsilon,\ldots,\psi^n_\epsilon$ be
orthonormal eigenforms of degree $p$ on $(\overline X_\epsilon,\bar
g_\epsilon)$ with eigenvalues tending to zero along a sequence
$\epsilon\to0$. Then, after passing to a subsequence, they converge
smoothly on compact subsets of $M_0$, and without loss of mass, in the sense that
$\langle\psi^a_\epsilon,\psi^b_\epsilon\rangle\to\langle\psi^a_0,
\psi^b_0\rangle$, to square-integrable closed and coclosed $p$-forms
$\psi^a_0$ for $g_0^R$. In particular the limits are orthonormal.
\end{proposition}
\begin{proof}
The model is of the kind treated by Ann\'e and Takahashi
\cite{anne-takahashi}: a fixed manifold with exactly conical ends at
finitely many points, to which complete manifolds with exactly conical
ends, scaled by $\epsilon$, are glued along the cones. Their analysis is carried out for one cap; nothing in it uses
connectedness of the cap, and it applies to finitely many, because the
tangential operator, the Poisson extensions and the boundary operators
on the caps, and the transmission conditions on the collars
\cite[Sections~2--4]{anne-takahashi} are direct sums over the nodes, and
the radial estimates are sums of integrals over disjoint collars. No
condition on the spectrum of the link is imposed there. Their
decomposition of a family of eigenforms with bounded eigenvalues
\cite[Corollary~17 and Lemma~18]{anne-takahashi} gives, on the central
part,
\[
 \psi^a_\epsilon=\Phi^a_\epsilon+\mathcal V^a_\epsilon
 +\mathcal C^a_\epsilon+o_{L^2}(1),\qquad
 \Phi^a_\epsilon\longrightarrow\psi^a_0\ \text{ in }L^2 ,
\]
where $\mathcal V^a_\epsilon$ comes from square-integrable solutions of
$(d+d^*)\psi=0$ on the complete caps and $\mathcal C^a_\epsilon$ from
solutions with $|\psi|=O(r^{-3})$, normalized by
$\ell_\epsilon^{-1/2}$, and the three terms are asymptotically
orthogonal:
\begin{equation}\label{eq:gram-limit}
 \lim\langle\psi^a_\epsilon,\psi^b_\epsilon\rangle
 =\langle\psi^a_0,\psi^b_0\rangle+\sum_i\langle v^a_i,v^b_i\rangle
 +\sum_i\langle\varsigma^a_i,\varsigma^b_i\rangle ,
\end{equation}
with $v_i^a$ the square-integrable cap solutions and $\varsigma_i^a$ the
coefficients of the solutions of order $r^{-3}$. By
Lemma~\ref{lem:cap-kernels}(1) the $v^a_i$ vanish. Since the eigenvalues
tend to zero, the limits $\psi^a_0$ are closed and coclosed of degree
$p$. The coefficients $\varsigma^a_i$ are leading coefficients of actual
solutions of $(d+d^*)\psi=0$ on the complete caps with
$|\psi|=O(r^{-3})$, by \cite[Proposition~11]{anne-takahashi} and the first
step of the proof of \cite[Proposition~12]{anne-takahashi}. By
Lemma~\ref{lem:cap-kernels}(2) the leading part of
$\mathcal C^a_\epsilon$ therefore has degree three; projecting the decomposition
onto forms of degree three, where $\psi^a_\epsilon$ and $\psi^a_0$ have
no component, shows that $\mathcal C^a_\epsilon\to0$ in $L^2$, so the
$\varsigma^a_i$ vanish. Then \eqref{eq:gram-limit} is the assertion.
\end{proof}

\begin{proposition}\label{prop:count}
The models have exactly $m$ positive eigenvalues in degree two tending to
zero as $\epsilon\to0$, and the degree-one spectrum and the remaining
positive degree-two spectrum are bounded below by a positive constant
depending on $R$, $A$ and $u$.
\end{proposition}
\begin{proof}
By Proposition~\ref{prop:compactness} and
Lemma~\ref{lem:central-kernels}, at most $b$ eigenvalues in degree two
and none in degree one tend to zero, counted with the zero eigenvalue.
Conversely, by Lemma~\ref{lem:central-kernels}(4) an orthonormal basis of
$\mathcal H_R$ is approximated in the graph norm of $d+d^*$ by compactly
supported forms; extended by zero to the model, they show by the min--max
principle that at least $b$ eigenvalues tend to zero. The model is
diffeomorphic to $X_t$ by Lemma~\ref{lem:condition-N}, so exactly
$s=b_2(X_t)$ of them are zero, by Lemma~\ref{lem:topology}.
\end{proof}

\subsection{The leading asymptotics}
Let $\bar\lambda_1(\epsilon)\le\cdots\le\bar\lambda_m(\epsilon)$ be the
small positive eigenvalues of the model in degree two.

\begin{theorem}\label{thm:model-law}
For fixed $R$, $A$ and $u$,
\[
 \lim_{\epsilon\to0}\log(1/\epsilon)\,\bar\lambda_j(\epsilon)=c\,\mu_j(R)
 \qquad(1\le j\le m),
\]
where $c=\pi$ and $\mu_j(R)$ are the eigenvalues of $\cE$ relative to
$\cM_R$.
\end{theorem}
\begin{proof}
\emph{The harmonic space of the model.} For $k\in K$ the form $\eta^R_k$
has no $\vartheta_i$-terms, so by \eqref{eq:cone-expansion} it equals
$d\beta_k$ near each node, and
$\gamma_{k,\rho}=\eta^R_k-d(\chi_\rho\beta_k)$, with $\chi_\rho$ equal to
one near the nodes and supported in $\{r<2\rho\}$, is a compactly
supported closed form in the class $k$ with
$\|\gamma_{k,\rho}-\eta^R_k\|^2=O(\rho^{2\kappa+2})$. Extended by zero it
is a closed form on the model representing the class in
$H^2(\overline X_\epsilon)$ that corresponds to $k$. Hence the harmonic
representative $h_\epsilon(k)$ of that class satisfies
$\limsup\|h_\epsilon(k)\|^2\le G_R(k,k)$. Conversely, by
Proposition~\ref{prop:compactness}, applied to the harmonic forms, a
subsequence converges without loss of norm to an element of
$\mathcal H_R$, whose class is $k$ by pairing with compactly supported
closed four-forms. Thus the Gram matrix of the harmonic forms of the
model on $K$ converges to $G_R|_K$, and the harmonic forms converge in
$L^2$ to the $\eta^R_k$.

\emph{Upper bound.} Let $D$ be $G_R$-orthogonal to $K$. With
\[
 \chi_\epsilon(r)=\frac{\log(r/2A\epsilon)}{\log(R/2A\epsilon)}
 \quad(2A\epsilon\le r\le R),
\]
extended by zero and one, put $T_\epsilon D=\chi_\epsilon\eta^R_D$,
extended by zero over the caps; it lies in the form domain of the model.
The metric is exactly conical where $\chi_\epsilon$ varies, and
$dT_\epsilon D=d\chi_\epsilon\wedge\eta^R_D$,
$d^*T_\epsilon D=-\iota_{\nabla\chi_\epsilon}\eta^R_D$. By
\eqref{eq:cone-expansion} and $\iota_{\partial_r}\vartheta_i=0$, the
leading term is $d\chi_\epsilon\wedge \per_i(D)\vartheta_i$, of squared norm
$c\,\per_i(D)^2\int r|\chi_\epsilon'|^2dr=c\,\per_i(D)^2/\log(R/2A\epsilon)$;
the cross term with $d\chi_\epsilon\wedge d\beta_D$ vanishes after
integration over the link because the tangential part of $d\beta_D$ is
exact on the link and $\vartheta_i$ is coclosed; and the remaining terms
are $O(\ell_\epsilon^{-2})\int_0^Rr^{2\kappa-1}dr$. Write $\bar q_\epsilon(\psi)=\|d\psi\|^2+\|d^*\psi\|^2$ for the energy
of a form on the model. Then
\begin{equation}\label{eq:trial}
\begin{split}
 \bar q_\epsilon(T_\epsilon D)&=\frac{c\,\cE([D],[D])}{\ell_\epsilon}
 +O(\ell_\epsilon^{-2})\|D\|^2,\\
 \langle T_\epsilon D,T_\epsilon E\rangle&=\cM_R([D],[E])
 +O(\ell_\epsilon^{-1})\|D\|\|E\| ;
\end{split}
\end{equation}
the second holds because $|1-\chi_\epsilon|\le C(1+|\log r|)/\ell_\epsilon$
is integrable against $|\eta^R_D|^2dV=O(r\,dr)$, and
because $G_R=\cM_R$ on the $G_R$-orthogonal complement of $K$. Subtract
from $T_\epsilon D$ its projection onto the harmonic forms of the model.
By the first step and $G_R(D,K)=0$ this changes the norm by $o(1)$, and
it does not change $\bar q_\epsilon$. The min--max principle on the span
of the first $j$ eigenvectors of $(\cE,\cM_R)$ gives
\[
 \limsup_{\epsilon\to0}\ell_\epsilon\bar\lambda_j(\epsilon)\le c\,\mu_j(R).
\]

\emph{Lower bound.} Choose a sequence along which
$\ell_\epsilon\bar\lambda_j(\epsilon)$ tends to its lower limit $\ell$.
By Proposition~\ref{prop:compactness} the first $j$ orthonormal
eigenforms $\psi^1_\epsilon,\ldots,\psi^j_\epsilon$ have, along a
subsequence, orthonormal limits $\psi^a_0\in\mathcal H_R$, which are
$L^2$-orthogonal to the $\eta^R_k$, $k\in K$, by the first step. For
$z\in\RR^j$ put $\psi_\epsilon(z)=\sum z_a\psi^a_\epsilon$ and
$\psi_0(z)=\sum z_a\psi^a_0=\eta^R_{D(z)}$, with $D(z)$ in the
$G_R$-orthogonal complement of $K$ and $\cM_R([D(z)])=|z|^2$. At the
fixed sections $\{r=R/2\}$ the forms converge smoothly, so
$\mathsf B_{\psi_\epsilon(z)}(R/2)\to \per_i(D(z))$ at the node $p_i$. The
part of the model inside this section at $p_i$ is isometric to a
truncated cap with the metric $\epsilon^2g^A_V$, by an identification
that differs from the chart by the unit $u$ and so preserves
$\vartheta_i$, and
Lemma~\ref{lem:cap-energy} gives
\[
 \bar\lambda_j(\epsilon)|z|^2\ge\bar q_\epsilon(\psi_\epsilon(z))\ge
 \sum_i\frac{c\,|\mathsf B_{\psi_\epsilon(z)}(R/2)|^2}
 {\log(R/2\epsilon)+C_A}.
\]
Multiplying by $\ell_\epsilon$ and passing to the limit,
$c\,\cE([D(z)])\le\ell\,\cM_R([D(z)])$ for all $z$. The classes $[D(z)]$
fill a $j$-dimensional subspace of $\cQ$, so $\ell\ge c\,\mu_j(R)$.
\end{proof}
